\chapter{Introduction}

Managed object-oriented languages such as Java free the programmer from manual memory
management: objects are created with \texttt{new} and are reclaimed automatically once they
are no longer reachable. This convenience has a price. Every object lives on a garbage-collected
heap, every access to it goes through a reference, and the runtime must periodically stop or
slow down the program to collect garbage. A long line of work has therefore tried to move
short-lived objects off the heap and onto the call stack, where they are cheap to allocate,
cheap to access, and disappear as soon as the method that created them returns.

This report studies that optimization, \emph{stack allocation}, in the context of the Eclipse
OpenJ9 Java virtual machine~\cite{openj9}. In this chapter we first describe how Java programs
are executed and how memory is managed, then explain how escape analysis enables stack
allocation, and finally discuss why getting a precise escape analysis is hard both inside a JIT
compiler and in a static compiler.

\section{Java Execution Pipeline}

Java programs are executed in two stages (Figure~\ref{fig:pipeline}). In the \emph{static}
stage, the \texttt{javac} compiler translates source code into platform-independent
bytecode, stored in \texttt{.class} files. In the \emph{dynamic} stage, the Java virtual machine
(JVM) loads this bytecode on the target machine and executes it. Every method starts in the
interpreter; methods that are invoked often enough are handed to the just-in-time (JIT) compiler,
which translates them into optimized machine code. From then on, calls to that method run the
compiled code instead of the interpreter, which is where most of the program's performance
comes from.

\begin{figure}[htbp]
  \centering
  \begin{tikzpicture}[node distance=1.1cm]
    \node[blk] (src) {Java source\\\texttt{.java}};
    \node[blk, right=1.3cm of src] (javac) {\texttt{javac}};
    \node[blk, right=1.3cm of javac] (bc) {Bytecode\\\texttt{.class}};
    \node[blk, right=1.6cm of bc, yshift=0.55cm, minimum width=2.7cm] (interp) {Interpreter};
    \node[blk, below=0.35cm of interp, minimum width=2.7cm] (jit) {JIT compiler(s)};
    \node[draw, dashed, rounded corners=3pt, inner sep=7pt, fit=(interp)(jit),
          label={[font=\small]above:JVM}] (jvm) {};
    \node[blk, right=1.0cm of jvm] (out) {Program\\output};
    \draw[arr] (src) -- (javac);
    \draw[arr] (javac) -- (bc);
    \draw[arr] (bc) -- (jvm.west);
    \draw[arr] (jvm) -- (out);
    \draw[densely dashed, gray] ($(bc.east)!0.5!(jvm.west) + (0,1.6)$) -- ++(0,-3.2);
    \node[font=\small\itshape, gray] at ($(javac) + (0,1.45)$) {Static};
    \node[font=\small\itshape, gray] at ($(jvm.north) + (0,0.85)$) {Dynamic};
  \end{tikzpicture}
  \caption{Java execution pipeline. \texttt{javac} produces bytecode statically; the JVM's
           interpreter and JIT compiler execute it dynamically.}
  \label{fig:pipeline}
\end{figure}

This split is central to the rest of the report. Analyses can be performed either statically,
where time is plentiful but the final program (its libraries, its dynamically loaded classes) is
not fully known, or during JIT compilation, where the running program is known exactly but every
millisecond spent analysing it is added to the program's execution time.

\section{Java Memory Management}

By default, Java allocates every object on the heap. The heap is owned and managed by the
garbage collector, which dictates how memory is structured and handles both the allocation and
the deallocation of objects.

\paragraph{Benefits.} Automatic memory management eliminates manual allocation and
deallocation, and with it a whole class of memory bugs such as dangling pointers, double frees
and many memory leaks.

\paragraph{Costs.} Every heap object is reached through a reference, so accesses pay for an
extra indirection, and the garbage collector adds its own performance overhead: it must trace
live objects, reclaim dead ones, and in many designs pause the application while doing so.

\section{Stack Allocation via Escape Analysis}

The JIT compiler can reduce these costs with \emph{escape analysis}.
Escape analysis determines the set of objects that do not \emph{escape} their allocating method,
that is, objects that are never stored into a location reachable after the method returns (a
static field, a heap object, the return value, or an argument of an unknown callee). Such
objects can be allocated on the method's stack frame instead of the heap. Stack-allocated
objects:
\begin{itemize}
  \item have a lower access time, since they live in the frame of the method using them; and
  \item are freed as soon as the allocating method returns, without any work from the garbage
        collector.
\end{itemize}

\paragraph{Limitations in OpenJ9.} OpenJ9's JIT compiler performs escape analysis and stack
allocation, but with two important limitations:
\begin{itemize}
  \item it is triggered only at the hotter optimization levels, so methods that are compiled
        only at lower levels never benefit from it; and
  \item it is intraprocedural, with a limited ability to ``peek'' into callees. Any call it
        cannot see through must be treated conservatively, which makes the analysis imprecise.
\end{itemize}

\section{Precision Challenges in Both Realms}
\label{sec:challenges}

A natural idea is to compute escape information more precisely. Doing so is hard in both stages
of the pipeline.

\paragraph{Challenges with JIT analysis.} The time spent in JIT compilation is added to the
execution time of the program. JIT compilers therefore trade precision for efficiency: they use
fast, local analyses rather than expensive whole-program ones.

\paragraph{Challenges with static analysis.} A static analysis can afford to be expensive, but it
faces two problems of its own:
\begin{itemize}
  \item a precise whole-program analysis needs the code of library methods, which may not be
        available (or may differ) when the application is compiled; and
  \item an object that the static analysis marks as stack-allocatable might still escape at run
        time, because of dynamic language features such as \emph{dynamic class loading},
        \emph{hot-code replacement}, and \emph{callbacks} into application code.
\end{itemize}

The work surveyed in Chapter~\ref{ch:literature} addresses both sets of challenges by splitting
the analysis between the static compiler and the JIT compiler.

\section{Organization of this Report}

Chapter~\ref{ch:literature} surveys two staged static\,+\,JIT approaches: the PYE framework and
optimistic stack allocation with dynamic heapification. Chapter~\ref{ch:toolbox} describes the
tools we use: the Soot analysis framework and OpenJ9's debug counters.
Chapter~\ref{ch:experiments} presents two experiments on OpenJ9, one on how well the stack-frame
budget for stack-allocated objects is used and one on why arrays are rarely stack allocated.
Chapter~\ref{ch:conclusion} concludes and outlines future research directions.
